\chapter{Introduction}
\label{ch:introduction}

This project exists to be compiled. If the PDF on the right shows four
chapters, a figure, a table and a bibliography, the whole pipeline works:
the engine ran, \texttt{bibtex} resolved the citations, and the document was
typeset more than once so the cross-references settled.

\section{Motivation}
\label{sec:motivation}

Typesetting was invented for mathematics~\cite{knuth1984texbook}, and the
document preparation system on top of it followed~\cite{lamport1994latex}.
Both entries live in \texttt{bibliography.bib}; typing \verb|\cite{| in the
editor should offer them by key, with the author and year alongside.

\section{What to look at}

Chapter~\ref{ch:methodology} contains a figure and a table.
Chapter~\ref{ch:results} contains the mathematics. Every
\verb|\ref| in this project resolves, so the Problems panel should be quiet
until you deliberately break something.

Inline maths sits in a sentence like $\vect{x} \in \mathbb{R}^{n}$, using the
\verb|\vect| command defined in \texttt{main.tex}.
